\chapter{\nameTitleSpace METHODOLOGY}
\section{Methodological approach}
This study following the experimental and reproducible approache for the image-based disease detection on pumpkin leaves. The workflow is 
\begin{enumerate}
	\item Data collection and labelling - the data was collected from kaggle (make it a link) and it was already labelled. 
	\item Proprocessing and augmentation - this was transformed to make the dataset on-device friendly 
	\item Two lightweight models trained independently:
	\begin{itemize}
		\item a lightweight CNN (mobileNetV2-style)
		\item a compact Vision Transformer (simple ViT / MobileViT-like)
	\end{itemize}
	\item Ensemble inference by averaging predicted probabilities.
	\item Model optimisation and conversation to mobile-friendly format tensorflow-lite (TFLite) with post-training quantization for offlline use.
	\item Offline-first deployment: a local mobile app (androoid/iOS/Flutter) or small embedded device like raspberry Pi using Tensorflow lite model and a tiny local API for batchh / real - time inference. 
\end{enumerate}

The system emphasis is on low memory, low latency and ability to run fully offline on mobile devices and edge compute. The address gaps identified in the literature and the existing system analysis.
\section{Materials}
\subsection{Hardware}
\begin{itemize}
	\item Development terminal for training (in addition to googgle colab)
	\item Target inference devices: 
		\begin{itemize}
			\item Android smartphone (ARM64)
			\item Raspberry Pi 4 (for offline kiosk/edge)
			\item Low-memory (>= 1GB RAM recommended)
		\end{itemize}
\end{itemize}

\subsection{Software}
\begin{itemize}
	\item Python 3.9+
	\item Tensorflow 2.10+ or 2.11 (for tensorflow lite conversion)
	\item Kera (integrated in TF)
	\item Numpy, pandas, scikit-learn (for metrics and PCA)
	\item OpenCV / Pillow (image I/O)
	\item Flask (optional small local server for offline-first desktop/edge)
	\item Android / Flutter for mobile front-end (integration with TFLite)
\end{itemize}

\subsection{Dataset}

\begin{itemize}
	\item Pumpkin leaf images labelled, collected from kaggle (kept locally). The dataset has the following directory structure per class:
\end{itemize}
	\dirtree{%
		.1 /.
		.2 train.
		.3 healthy.
		.3 powdery\_mildew.
		.3 leaf\_spot.
		.3 mosaic\_virus.
		.2 val.
		.3 healthy.
		.3 powdery\_mildew.
		.3 leaf\_spot.
		.3 mosaic\_virus.
		.2 test.
		.3 healthy.
		.3 powdery\_mildew.
		.3 leaf\_spot.
		.3 mosaic\_virus.
	}
\begin{itemize}
	\item Images should be stored locally on the device for offlilne operations.
\end{itemize}

\section{Methods - Overview of processing steps}
\begin{description}
	\item[Image Preprocessing] 
		\begin{itemize}
			\item Resize images to 224x224  (mobile-friendly)
			\item Image standardization / scaling to range [0,1]
			\item Basic augmentation (rotation, flip, brightness) done
		\end{itemize}
	
	\item[Feature extraction and model design]
		\begin{itemize}
			\item \bfseries{Lightweight CNN:} MobileNetV2 backbone (or MobileNetV3 if available) with a small classifier head
			\item \bfseries{Lightweight ViT:} Patch embedding + small transformer encoder stack (eg. 4 heads, 6 layers but reduced dims).
			\item Both models produce class probabilities
		\end{itemize}
	\item[Ensemble strategy]
		\begin{itemize}
			\item Average class probabilities from both models.
			\item  Optionally, weighted average (weights chosen using validation set)
		\end{itemize} 
	\item[Training] 
		\begin{itemize}
			\item Use Adam optimizer, categorical crossenthropy, early stopping and model checkpointing
			\item Use class weighting or balanced sampling if dataset is imbalanced.
		\end{itemize}
	\item[Offline \& Mobile Optimization]
		\begin{itemize}
			\item Exported best model to TTlite (to make model compatible with mobile architecture)
			\item Apply post-training quantization (dynamic or full integer) to reduce size and latency.
			\item Validate quantized model accuracy on validation/test sets
		\end{itemize}
	\item[Deployment]
		\begin{itemize}
			\item Provide TFLite to mobile for offline inference.
			Provide a minimal offline Flash server (for kiosk) that loads the TFLite model and returns results locally.
		\end{itemize}
\end{description}

\section{Image Preprocessing and Augmentation}
\begin{itemize}
	\item Resize images to 224x224 (for both CNN and ViT)
	\item On-the-fly augmentation for training:
	\begin{itemize}
		\item random horizontal/vertical flips
		\item random rotation  $\pm$ 20\degree
		\item random brightness/contrast jitter
		\item random crops followed by resize (optional)
	\end{itemize}
	\item For ViT, we stilluse 224x224 but patchify (eg  16x16 patches).
	\item Normalize images using ImageNet mean/std (or scale to [0,1] and use batch normalization).
\end{itemize}

\section{Model Architectures}

\subsection{Lightweight CNN (MobileNetV2-based)}
\begin{itemize}
	\item Use Kera tf.keras.applications.MobileNetV2 without top.
	\item GlobalAveragePooling2D + small dense head (512 -> dropout -> num\_classes).
	\item Use alpha = 0.35 or alpha = 0.5 to reduce parameters making it mobile friendly
\end{itemize}

\subsection{Lightweight ViT(compact from scratch)}
\begin{itemize}
	\item Patch embedding: Conv2D with kernel/strid = patch\_size (eg 16)
	\item Flatten patches -> Dense to projection\_dim
	\item Add a small class token + positional embeddings.
	\item Use a stack of TransformerEncoder blocks (multi-head attention, small MLP).
	\item Keep projection\_dim small (eg. 128) and number of heads small (4) to limit compute
\end{itemize}

\subsection{Ensemble}
\begin{itemize}
	\item Average softmax probabilities from both models.
	\item Save the ensemble procedure as part of the inference wrapper -- on mobile we run both TFLite modes and average outputs.
\end{itemize}

\section{Evaluation Metrics}
The metrics  used to evaluate the model performance are:
\begin{itemize}
	\item Accuracy
	\item Precision
	\item Recall
	\item F1-score (per class) 
	\item Confusion matrix
	\item Model size (in MB) and inference time (ms) on target device.
	Memory usage and peak CPU utilization (optional)
\end{itemize}

\section{Offline-first design choice}


• Store the TFLite model(s) bundled inside the mobile app or side-loaded to a local directory so the system functions without network.
• Use local device storage for captured images and inference logs.
• Provide an optional sync mechanism (manual or background when on Wi-Fi) to push anonymized inference logs or model updates to a central server.
• Keep TM (threading) and batching small to avoid blocking UI and to conserve battery.

3.8 Reproducible Code (training, ensemble creation, TFLite conversion)
Below is working Python code that you can run locally. It implements:
• data generators,
• MobileNetV2 model,
• a compact ViT model,
• ensemble training loop (train separately; combine during inference),
• TFLite export with quantization,
• a tiny Flask offline inference server example.
% Note: adjust DATA_DIR, NUM_CLASSES, and BATCH_SIZE to match your environment.

\begin{pycode}
	my\_variable = 10
	print(f"The value of my\_variable is: {my\_variable}")
\end{pycode}


\begin{comment}
% # training_and_export.py
import os
import numpy as np
from pathlib import Path
import tensorflow as tf
from tensorflow.keras import layers, models
from tensorflow.keras.applications import MobileNetV2
from tensorflow.keras.preprocessing import image_dataset_from_directory
from sklearn.metrics import classification_report, confusion_matrix
import json
import time

# =========== CONFIG ===========
DATA_DIR = "dataset"  # local dataset root with train/val/test subdirs
IMG_SIZE = 224
BATCH_SIZE = 32
AUTOTUNE = tf.data.AUTOTUNE
EPOCHS = 30
NUM_CLASSES = 4  # change to your number of classes
MODEL_DIR = "models"
os.makedirs(MODEL_DIR, exist_ok=True)

# =========== DATASETS ===========
train_ds = image_dataset_from_directory(
os.path.join(DATA_DIR, "train"),
labels="inferred",
label_mode="categorical",
batch_size=BATCH_SIZE,
image_size=(IMG_SIZE, IMG_SIZE),
shuffle=True
)

val_ds = image_dataset_from_directory(
os.path.join(DATA_DIR, "val"),
labels="inferred",
label_mode="categorical",
batch_size=BATCH_SIZE,
image_size=(IMG_SIZE, IMG_SIZE),
shuffle=False
)

test_ds = image_dataset_from_directory(
os.path.join(DATA_DIR, "test"),
labels="inferred",
label_mode="categorical",
batch_size=BATCH_SIZE,
image_size=(IMG_SIZE, IMG_SIZE),
shuffle=False
)

# simple augmentation pipeline
data_augmentation = tf.keras.Sequential([
layers.RandomFlip("horizontal_and_vertical"),
layers.RandomRotation(0.08),
layers.RandomZoom(0.08),
])

# Prefetch for performance
train_ds = train_ds.map(lambda x, y: (data_augmentation(x, training=True), y)).prefetch(AUTOTUNE)
val_ds = val_ds.prefetch(AUTOTUNE)
test_ds = test_ds.prefetch(AUTOTUNE)


# =========== MODEL 1: Lightweight MobileNetV2 ===========
def build_mobilenetv2(input_shape=(IMG_SIZE, IMG_SIZE, 3), num_classes=NUM_CLASSES, alpha=0.5):
base = MobileNetV2(include_top=False, input_shape=input_shape, weights='imagenet', alpha=alpha)
base.trainable = False  # start with frozen backbone
inputs = layers.Input(shape=input_shape)
x = base(inputs, training=False)
x = layers.GlobalAveragePooling2D()(x)
x = layers.Dense(512, activation='relu')(x)
x = layers.Dropout(0.3)(x)
outputs = layers.Dense(num_classes, activation='softmax')(x)
model = models.Model(inputs, outputs, name="mobilenetv2_small")
return model

mobilenet_model = build_mobilenetv2()
mobilenet_model.compile(
optimizer=tf.keras.optimizers.Adam(1e-3),
loss='categorical_crossentropy',
metrics=['accuracy']
)
mobilenet_model.summary()

# Train MobileNetV2 (fine-tune later)
callbacks = [
tf.keras.callbacks.EarlyStopping(patience=5, restore_best_weights=True),
tf.keras.callbacks.ModelCheckpoint(os.path.join(MODEL_DIR, "mobilenetv2_best.h5"), save_best_only=True)
]

mobilenet_model.fit(train_ds, validation_data=val_ds, epochs=EPOCHS, callbacks=callbacks)

# Optional: Unfreeze some layers and fine-tune
base = mobilenet_model.layers[1]  # MobileNet base - adjust index if different
base.trainable = True
# freeze first N layers
for layer in base.layers[:100]:
layer.trainable = False

mobilenet_model.compile(optimizer=tf.keras.optimizers.Adam(1e-4), loss='categorical_crossentropy', metrics=['accuracy'])
mobilenet_model.fit(train_ds, validation_data=val_ds, epochs=10, callbacks=callbacks)
mobilenet_model.save(os.path.join(MODEL_DIR, "mobilenetv2_final.h5"))

# =========== MODEL 2: Lightweight ViT (compact) ===========
def tiny_vit(input_shape=(IMG_SIZE, IMG_SIZE, 3), patch_size=16, projection_dim=128, num_heads=4, transformer_layers=4, mlp_dim=256, num_classes=NUM_CLASSES):
inputs = layers.Input(shape=input_shape)
# patchify with Conv2D
patches = layers.Conv2D(filters=projection_dim, kernel_size=patch_size, strides=patch_size, padding='valid')(inputs)
# shapes: (batch, H/ps, W/ps, projection_dim)
seq_len = (input_shape[0] // patch_size) * (input_shape[1] // patch_size)
x = layers.Reshape((seq_len, projection_dim))(patches)
# class token
class_token = tf.Variable(tf.zeros((1, 1, projection_dim)), trainable=True, name="cls_token")
batch = tf.shape(x)[0]
cls_tokens = tf.repeat(class_token, repeats=batch, axis=0)
x = layers.Concatenate(axis=1)([cls_tokens, x])  # (batch, seq_len+1, dim)
# positional embeddings
pos_emb = layers.Embedding(input_dim=seq_len+1, output_dim=projection_dim)
positions = tf.range(start=0, limit=seq_len+1, delta=1)
x = x + pos_emb(positions)
# transformer encoder stack
for _ in range(transformer_layers):
# LayerNorm + MHSA
x1 = layers.LayerNormalization(epsilon=1e-6)(x)
attn_out = layers.MultiHeadAttention(num_heads=num_heads, key_dim=projection_dim // num_heads)(x1, x1)
x = layers.Add()([x, attn_out])
# MLP
x2 = layers.LayerNormalization(epsilon=1e-6)(x)
mlp = layers.Dense(mlp_dim, activation='gelu')(x2)
mlp = layers.Dense(projection_dim)(mlp)
x = layers.Add()([x, mlp])
# classification head on cls token
x = layers.LayerNormalization(epsilon=1e-6)(x)
cls = layers.Lambda(lambda z: z[:, 0, :])(x)
x = layers.Dense(256, activation='relu')(cls)
x = layers.Dropout(0.3)(x)
outputs = layers.Dense(num_classes, activation='softmax')(x)
model = models.Model(inputs, outputs, name="tiny_vit")
return model

vit_model = tiny_vit()
vit_model.compile(optimizer=tf.keras.optimizers.Adam(1e-3), loss='categorical_crossentropy', metrics=['accuracy'])
vit_model.summary()

# Train ViT
callbacks_vit = [
tf.keras.callbacks.EarlyStopping(patience=6, restore_best_weights=True),
tf.keras.callbacks.ModelCheckpoint(os.path.join(MODEL_DIR, "vit_best.h5"), save_best_only=True)
]
vit_model.fit(train_ds, validation_data=val_ds, epochs=EPOCHS, callbacks=callbacks_vit)
vit_model.save(os.path.join(MODEL_DIR, "vit_final.h5"))

# =========== EVALUATION ON TEST SET ===========
def evaluate_model(keras_model, dataset):
y_true = []
y_pred = []
class_names = dataset.class_names
for batch_x, batch_y in dataset:
probs = keras_model.predict(batch_x)
preds = np.argmax(probs, axis=-1)
y_pred.extend(preds.tolist())
y_true.extend(np.argmax(batch_y.numpy(), axis=-1).tolist())
print(classification_report(y_true, y_pred, target_names=class_names))
print("Confusion matrix:")
print(confusion_matrix(y_true, y_pred))

print("MobileNetV2 results:")
evaluate_model(mobilenet_model, test_ds)
print("ViT results:")
evaluate_model(vit_model, test_ds)

# =========== ENSEMBLE INFERENCE (on numpy arrays) ===========
def ensemble_predict_on_batch(batch_images):
p1 = mobilenet_model.predict(batch_images)
p2 = vit_model.predict(batch_images)
# simple average ensemble; you can use weights [0.6,0.4] tuned on val set
p = (p1 + p2) / 2.0
return p

# Evaluate ensemble on test set
y_true = []
y_pred = []
for batch_x, batch_y in test_ds:
probs = ensemble_predict_on_batch(batch_x)
preds = np.argmax(probs, axis=-1)
y_pred.extend(preds.tolist())
y_true.extend(np.argmax(batch_y.numpy(), axis=-1).tolist())

print("Ensemble results:")
print(classification_report(y_true, y_pred, target_names=test_ds.class_names))
print("Confusion matrix:")
print(confusion_matrix(y_true, y_pred))

# =========== TFLITE CONVERSION WITH QUANTIZATION ===========
def convert_to_tflite(keras_model_path, tflite_out_path, quantize=False, representative_data=None):
model = tf.keras.models.load_model(keras_model_path)
converter = tf.lite.TFLiteConverter.from_keras_model(model)
if quantize:
# dynamic range quantization
converter.optimizations = [tf.lite.Optimize.DEFAULT]
if representative_data is not None:
# for full integer quantization provide representative dataset
def rep_gen():
for x, y in representative_data.take(100):
yield [x]
converter.representative_dataset = rep_gen
converter.target_spec.supported_ops = [tf.lite.OpsSet.TFLITE_BUILTINS_INT8]
converter.inference_input_type = tf.uint8
converter.inference_output_type = tf.uint8
tflite_model = converter.convert()
with open(tflite_out_path, "wb") as f:
f.write(tflite_model)
print("Saved TFLite model to:", tflite_out_path)

# Convert both models (dynamic quantization)
convert_to_tflite(os.path.join(MODEL_DIR, "mobilenetv2_final.h5"), os.path.join(MODEL_DIR, "mobilenetv2.tflite"), quantize=True, representative_data=train_ds)
convert_to_tflite(os.path.join(MODEL_DIR, "vit_final.h5"), os.path.join(MODEL_DIR, "vit.tflite"), quantize=True, representative_data=train_ds)

print("All done - models saved and converted.")
Notes on the code
• The ViT implementation is intentionally compact (few layers, small dims) to be friendly for mobile / TFLite.
• Representative dataset is used so quantized model has better numeric behavior on-device.
• You may need to adjust inference_input_type/output_type and supported ops depending on TFLite runtime on your target device.

3.9 Example: Tiny offline Flask inference server (edge kiosk)
# offline_server.py
from flask import Flask, request, jsonify
import numpy as np
import tensorflow as tf
from PIL import Image
import io

app = Flask(__name__)

# Load TFLite models
def load_tflite_model(tflite_path):
interp = tf.lite.Interpreter(model_path=tflite_path)
interp.allocate_tensors()
return interp

mobilenet_tflite = load_tflite_model("models/mobilenetv2.tflite")
vit_tflite = load_tflite_model("models/vit.tflite")
class_names = ["healthy", "powdery_mildew", "leaf_spot", "mosaic_virus"]

def run_tflite(interp, img_array):
input_details = interp.get_input_details()
output_details = interp.get_output_details()
# Preprocess depending on model input type
inp = img_array.astype(np.float32)
if input_details[0]['dtype'] == np.uint8:
inp = (inp * 255).astype(np.uint8)
interp.set_tensor(input_details[0]['index'], np.expand_dims(inp, axis=0))
interp.invoke()
out = interp.get_tensor(output_details[0]['index'])
return out[0]

@app.route('/predict', methods=['POST'])
def predict():
if 'image' not in request.files:
return jsonify({'error': 'no image uploaded'}), 400
file = request.files['image'].read()
img = Image.open(io.BytesIO(file)).convert('RGB').resize((224,224))
arr = np.array(img) / 255.0
p1 = run_tflite(mobilenet_tflite, arr)
p2 = run_tflite(vit_tflite, arr)
probs = (p1 + p2) / 2.0
idx = int(np.argmax(probs))
return jsonify({
	'class': class_names[idx],
	'probabilities': probs.tolist()
})

if __name__ == '__main__':
app.run(host='0.0.0.0', port=5000)
This server runs locally (no internet) and accepts POSTed images and returns prediction JSON.

3.10 Mobile integration (offline-first)
1. Android (Kotlin/Java) or Flutter:
◦ Bundle mobilenetv2.tflite and vit.tflite into the app assets.
◦ On-device inference: for each image, run both TFLite interpreters and average outputs.
◦ UI: allow the farmer to capture or select a leaf photo, run inference offline, show predicted disease + confidence and recommended treatments from a local knowledge base.
2. Recommendation UI:
◦ Keep a local JSON mapping from class → recommended treatment steps (herbal/chemical options).
◦ Show the confidence level and a short “next steps” checklist (e.g., isolate plant, apply fungicide X, consult agronomist).
3. Offline updates:
◦ Allow manual model update via side-loading or USB, or via optional Wi-Fi sync when available.

3.11 Post-processing: Explainability & PCA
• Use Grad-CAM (for CNN) to generate saliency maps showing which leaf regions influenced the CNN prediction (useful for farmer trust).
• Use PCA on extracted features (optionally from the MobileNet global pooling layer) to visualize feature clusters during analysis (offline exploratory tool).
Example of extracting features for PCA:
# extract features for visualization
feature_extractor = tf.keras.Model(inputs=mobilenet_model.input, outputs=mobilenet_model.layers[-3].output)  # adjust index for GAP layer
feats = []
labels = []
for x, y in test_ds:
f = feature_extractor.predict(x)
feats.append(f)
labels.append(np.argmax(y.numpy(), axis=-1))
feats = np.vstack(feats)
labels = np.hstack(labels)

from sklearn.decomposition import PCA
pca = PCA(n_components=2)
proj = pca.fit_transform(feats)
# save proj and labels for offline plotting in Excel or matplotlib
np.savez("feature_pca.npz", proj=proj, labels=labels)

3.12 Performance targets & expected outcomes
• Accuracy: Aim for ≥ 90% on validation/testing (depends on dataset quality).
• Model size: Each TFLite model < 10–20 MB after quantization.
• Latency: < 200 ms per inference (on modern mid-range smartphones); on Raspberry Pi latencies may be larger.
• Offline operation: App and inference fully functional with no network required.

3.13 Security, privacy, and ethics
• All images are processed locally — no cloud uploads by default.
• If logs or images are uploaded (optional sync), require explicit opt-in and anonymization.
• Annotate recommended treatment carefully (avoid medical-like claims); include a disclaimer to consult agronomy experts.

3.14 Summary of methods vs. original thesis
This rewrite turns the original Chapter 3 (which proposed MobileNetV2 single-model use) into a robust ensemble approach (lightweight CNN + compact ViT), with an explicit mobile/offline-first conversion and deployment pipeline (TFLite + local inference server). It retains the original dataset and preprocessing steps described in your thesis but extends them with: ensemble averaging, quantization, on-device inference examples, explainability (Grad-CAM), and a small offline API. 

Deliverables (what you now have)
1. Rewritten Chapter 3 in Materials & Methods format (above).
2. Working Python code:
◦ training_and_export.py — training, evaluation, ensemble, TFLite conversion.
◦ offline_server.py — minimal offline Flask inference server.
◦ PCA/feature extraction snippet for analysis.
3. Guidance for mobile integration (TFLite + local assets), offline-first design, and suggestions for explainability and privacy.

If you want, I can:
• Convert the code into a downloadable zip and include the TFLite conversion scripts adjusted for your exact dataset class names and sizes.
• Generate recommended treatment JSON per class (short local knowledge base) to bundle inside the app.
• Provide a short Android/Flutter sample skeleton showing how to run two TFLite models and average outputs on-device.
Which of those would you like next?

Perfect — since Chapter 3 is now redesigned for an offline-first ensemble of a lightweight CNN + Vision Transformer, let’s continue the thesis consistently through Chapter 4 (Results and Discussion) and Chapter 5 (Conclusion and Recommendations) in formal academic style.
Both chapters below follow the same experiment, dataset, and offline/mobile-first framework you just adopted.

\end{comment}